% chktex-file 1 
% chktex-file 2
% chktex-file 3
% chktex-file 9
% chktex-file 11
% chktex-file 17
% chktex-file 36

%%%%GPT writing assistance - the fundamental proof idea comes from the author - GPT Astra identified the following flaw: the Hilbertspaces H_L and \H^{-\sigma} needed be equivalent for the proof of theorem \ref{theorem:uniqueness_SDMFW} to be true and it provided the Lemmas resolving this problem.
%%%%Furthemore, it streamlined proofs considerably 
\section{Uniqueness of the invariant dynamics of the SDMFW}\label{Chpt_Invariant_Dynamics}

Let $\sigma\in(0,\frac12)$ and let
\[
    \mu\in\mathcal P\bigl(\mathcal H^{-\sigma}(\T^2)\bigr)
    \qquad \text{with} \qquad
    \mathcal H^{-\sigma}
    =H^{-\sigma}(\T^2)\times H^{-1-\sigma}(\T^2),
\]
be an invariant law for \eqref{eq:SDMFW} such that
\[
    M_\mu:=\mathbb E_\mu[\wick{u^2}]
    \in H^{-\sigma}(\T^2).
\]
As explained in Section \ref{sec:SDMFW_inv_msr}, stationarity freezes
the mean-field term.  Hence, a stationary solution satisfies the linear
equation
\begin{equation}\label{eq:frozen_linear_dynamics}
    \begin{cases}
        du_\eps=v_\eps\,dt,\\
        dv_\eps=-\eps^{-2}(L_\mu u_\eps+v_\eps)\,dt
        +\sqrt2\,\eps^{-2}\,dW_t,
    \end{cases}
\end{equation}
with $\Lambda:=1-\Delta$ and $L_\mu:=\Lambda+M_\mu$. 

\noindent
We first construct $L_\mu$ as a self-adjoint operator with compact
resolvent that is bounded from below. We then prove sufficient Sobolev regularity and show that $L_\mu$ induces a Hilbert Space $\H_L$ with norm equivalent to the canonical norm of $\H^{-\sigma}$. Building on the framework established in \cite{DaPrato_Zabczyk_2014} to analyze stochastic differential equations (SDEs) in infinite dimensions, we prove that~\eqref{eq:frozen_linear_dynamics} admits a weak solution in Proposition~\ref{prop:mild} (in the sense of Theorem 5.4 in \cite{DaPrato_Zabczyk_2014}). Then, using the characterization Theorem 11.7 in~\cite{DaPrato_Zabczyk_2014} of invariant measures of linear infinite dimensional SDEs, we uniquely determine the invariant measure $\mu = \mathcal{N}(0, L_\mu^{-1}) \otimes \mathcal{N}(0, \varepsilon^{-2}I)$ of~\eqref{eq:frozen_linear_dynamics}. 
Finally, we are left to show that self-consistency forces $M_\mu = 0$ which implies that any invariant measure $\mu$ of the~\eqref{eq:SDMFW} equals $\rho_\eps = \mathcal{N}(0, \Lambda^{-1}) \otimes \mathcal{N}(0, \varepsilon^{-2}I)$.

\subsection{The operator \(L_{\mu}\) and its spectral basis}\label{sec_invdynamics_mildsolution}

For the remainder of this section we write $M = M_{\mu}$ and $L = L_{\mu}$.
The real-valued distribution $M$ defines on $H^1(\T^2)$ the symmetric
quadratic form
\[
    \mathfrak l[u,v]
    :=\langle \Lambda^{1/2}u,\Lambda^{1/2}v\rangle_{L^2_x}
      +\langle M,u\overline v\rangle,
\]
where the second term is understood by Sobolev duality.

This first lemma utilizes the KLMN Theorem to prove that the operator $L$ can be realized as a self-adjoint operator on $L^2(\T^2)$.
\begin{lemma}\label{lem:relative_bound}
Suppose $M\in H^{-\sigma}(\T^2)$ for some $\sigma\in(0,1)$.
For every $\delta>0$ there exists $C_\delta>0$ such that
\[
    |\langle M,|u|^2\rangle|
    \leq
    \delta\langle\Lambda u,u\rangle_{L^2_x}
    +C_\delta\|u\|_{L^2_x}^2,
    \qquad u\in H^1(\T^2).
\]
Consequently, $\mathfrak l$ is closed and bounded from below on
$H^1(\T^2)$ and determines a self-adjoint operator $L$ on
$L^2(\T^2)$ with form domain $H^1(\T^2)$.
\end{lemma}

\begin{proof}
Choose $\eta>0$ such that $2\eta+\sigma\leq1$.  By Sobolev duality,
the negative-order product estimate in \cite[Proposition~2]{Benyi2025},
and Sobolev embedding,
\begin{align*}
    |\langle M,|u|^2\rangle|
    \leq
    \|Mu\|_{H^{-1+\eta}_x}\|u\|_{H^{1-\eta}_x}
    \lesssim
    \|M\|_{H^{-\sigma}_x}\|u\|_{H^{1-\eta}_x}^2.
\end{align*}
Interpolating between fractional Sobolev spaces, for any $\tilde{\delta} > 0$, there exists a constant $C(\tilde{\delta}) > 0$ such that
\[
    \|u\|_{H^{1-\eta}_x}^2
    \leq
    \widetilde\delta\|u\|_{H^1_x}^2
    +C(\widetilde{\delta})\|u\|_{L^2_x}^2.
\]
Taking $\widetilde\delta$ sufficiently small proves the stated bound.
The final assertion follows from the KLMN theorem
\cite[Theorem~10.21]{schmudgen_operators}.
\end{proof}

The next lemma demonstrates that $L$ has compact resolvent and characterizes its spectrum in relation to the operator $\Lambda$. 
\begin{lemma}\label{lemma_equ_ev}
Suppose $M\in H^{-\sigma}(\T^2)$ for some $\sigma\in(0,1)$. There exists $c>0$ such that $L_c:=L+c$ is strictly positive and
\begin{equation}\label{eq:form_equivalence}
    C_1\langle\Lambda u,u\rangle_{L^2_x}
    \leq
    \langle L_cu,u\rangle_{L^2_x}
    \leq
    C_2\langle\Lambda u,u\rangle_{L^2_x},
    \qquad u\in H^1(\T^2),
\end{equation}
for some $C_1,C_2>0$.  Moreover, $L$ has compact resolvent.  In
particular, there exist a real orthonormal basis $\{e_k\}_{k\in \N}$ of
$L^2(\T^2)$ and real eigenvalues $(\lambda_k)_{k\geq1}$, repeated
according to multiplicity and satisfying $\lambda_k\to+\infty$, such
that
\[
    Le_k=\lambda_ke_k.
\]
If $(\nu_k)_{k\geq1}$ are the eigenvalues of $L_c$ and
$(\ell_k)_{k\geq1}$ those of $\Lambda$, both in nondecreasing order,
then
\[
    C_1\ell_k\leq\nu_k\leq C_2\ell_k,
    \qquad k\geq1.
\]
\end{lemma}

\begin{proof}
By Lemma \ref{lem:relative_bound}, $L$ is bounded from below.  Choose
$c>0$ so that $L_c$ is strictly positive.  The relative form bound,
combined with strict positivity of $L_c$, gives
\eqref{eq:form_equivalence}.

For $f\in L^2(\T^2)$ and $u=L_c^{-1}f$, the lower bound in
\eqref{eq:form_equivalence} gives
\[
    \|u\|_{H^1_x}^2
    \lesssim
    \langle L_cu,u\rangle_{L^2_x}
    =\langle f,u\rangle_{L^2_x}
    \lesssim
    \|f\|_{L^2_x}\|u\|_{H^1_x}.
\]
Thus, $L_c^{-1}:L^2(\T^2)\to H^1(\T^2)$ is bounded.  Since
$H^1(\T^2)\hookrightarrow L^2(\T^2)$ is compact, $L_c^{-1}$ is
compact on $L^2$; hence both $L_c$ and $L$ have compact resolvent.
The spectral assertions follow from the spectral theorem.  Finally,
the eigenvalue comparison follows from \eqref{eq:form_equivalence} and
the min-max principle.
\end{proof}
Since the previous Lemma establishes that for some $c > 0$, $L_c \coloneq L + c$ is a strictly positive self-adjoint operator, we can define $\H_L$ as the closure of $L^2(\T^2) \times L^2(\T^2)$ with respect to the norm induced by the inner product
\begin{align}\label{eq:inner_product}
    \langle \mathbf{u}, \mathbf{v} \rangle_{L_c} \coloneqq \langle L_c^{-\sigma} u_1, v_1 \rangle_{L^2_x} + \langle L_c^{-1-\sigma} u_2, v_2 \rangle_{L^2_x},
\end{align}
with respect to the normalized inner product on $L^2(\T^2)$:
\begin{align*}
    \langle u, v \rangle_{L^2_x} \coloneq (2\pi)^{-2} \int_{\T^2} u(x) \overline{v(x)} \, dx.
\end{align*}

Finally, the following lemma is the point at which the restriction $\sigma<\frac12$ enters the spectral projection argument. The Corollary following this Lemma, infers that the Hilbert spaces $\H_L$ and $\H^{-\sigma}$ are equivalent. 
%%%%GPT Astra very high: this Lemma was conceived by Chat GPT itself - the author only added details. 
\begin{lemma}\label{lem:eigenfunction_regularity}
Let $\sigma\in(0,\frac12)$ and $M\in H^{-\sigma}(\T^2)$.  For the
constant $c$ in Lemma \ref{lemma_equ_ev}, the operator $L_c$ extends
to an isomorphism
\begin{equation}\label{eq:Lc_sobolev_isomorphism}
    L_c:H^{1+\sigma}(\T^2)
    \longrightarrow H^{-1+\sigma}(\T^2).
\end{equation}
Consequently,
\[
    e_k\in H^{1+\sigma}(\T^2),
    \qquad k\geq1,
\]
and the algebraic span
\[
    \mathscr E_L:=\operatorname{span}\{e_k:k\geq1\}
\]
is dense in both $H^{1+\sigma}(\T^2)$ and
$H^\sigma(\T^2)$.
\end{lemma}

\begin{proof}
From the product estimate \cite{Benyi2025} Proposition 2 and Sobolev's inequality if follows that multiplication by an element of
$H^{1+\sigma}(\T^2)$ is bounded on $H^{-\sigma}(\T^2)$.  Hence,
\begin{equation}\label{eq:potential_multiplier_bound}
    \|Mf\|_{H^{-\sigma}_x}
    \lesssim
    \|M\|_{H^{-\sigma}_x}\|f\|_{H^{1+\sigma}_x}.
\end{equation}
Because $0 < \sigma<\frac12$, the embedding
\[
    H^{-\sigma}(\T^2)
    \hookrightarrow H^{-1+\sigma}(\T^2)
\]
is compact.  It follows that multiplication by $M$ is a compact
operator from $H^{1+\sigma}(\T^2)$ to $H^{-1+\sigma}(\T^2)$.  The same is true of
the zeroth-order term $cI$.  Since
\[
    \Lambda:H^{1+\sigma}(\T^2)
    \longrightarrow H^{-1+\sigma}(\T^2)
\]
is an isomorphism, $K = \Lambda^{-1}(M+cI)$ is compact. Additionally, $L_c$ has trivial kernel: if $L_cf=0$, then
$f\in H^{1+\sigma}(\T^2)\subset H^1(\T^2)$, and strict positivity of the form of $L_c$ gives $f=0$.
Thus, writing $L_c=\Lambda+M+c = \Lambda(I + K)$, the Fredholm alternative \cite[Theorem~6.6]{Brezis2011} implies that $L_c$ is Fredholm of index zero between
these spaces. Finally, from the bounded inverse theorem follows \eqref{eq:Lc_sobolev_isomorphism}.

Since $(\lambda_k+c)e_k\in L^2(\T^2)\subset H^{-1+\sigma}(\T^2)$,
\eqref{eq:Lc_sobolev_isomorphism} gives $e_k\in H^{1+\sigma}(\T^2)$.
It remains to prove density.  Let $f\in H^{1+\sigma}(\T^2)$ and set
$g=L_cf\in H^{-1+\sigma}(\T^2)$.  The finite linear span of $\{e_k\}_{k \in \N}$ is
dense in $H^{-1+\sigma}(\T^2)$: indeed, $L^2(\T^2)$ is dense in
$H^{-1+\sigma}(\T^2)$, while the finite eigenfunction expansions are dense
in $L^2(\T^2)$.  Choose $g_N\in\mathscr E_L$ with
$g_N\to g$ in $H^{-1+\sigma}(\T^2)$ and set $f_N=L_c^{-1}g_N$.  Then
$f_N\in\mathscr E_L$ and, by \eqref{eq:Lc_sobolev_isomorphism},
$f_N\to f$ in $H^{1+\sigma}(\T^2)$.  Thus, $\mathscr E_L$ is dense in
$H^{1+\sigma}(\T^2)$.  Its density in $H^\sigma(\T^2)$ follows from the density
of $H^{1+\sigma}(\T^2)$ in $H^\sigma(\T^2)$.
\end{proof}
For \(s\geq0\), equip
\[
 H_{L_c}^s:=D(L_c^{s/2})
\]
with the norm 
\[
\|f\|_{H_{L_c}^s}:=\|L_c^{s/2}f\|_{L^2_x}.
\]
For \(s<0\), let \(H_{L_c}^s\) be the completion of \(L^2(\T^2)\)
in this norm, equivalently the \(L^2\)-pivot dual of \(H_{L_c}^{-s}\).

\begin{corollary}[Identification of the Sobolev phase space]\label{cor:equ_norms}
For every \(s\in[-1-\sigma,1+\sigma]\),
\begin{align*}
    H_{L_c}^s=H^s(\T^2)
\end{align*}
with equivalent norms. In particular,
\begin{align*}
 \H_L=H_{L_c}^{-\sigma}\times H_{L_c}^{-1-\sigma}
     =\H^{-\sigma}.
\end{align*}
These identifications preserve the Borel sigma-algebras.
\end{corollary}

\begin{proof}

Form equivalence of $L_c$ and $\Lambda$, established in Lemma~\ref{lemma_equ_ev}, identifies \(H_{L_c}^1\) with \(H^1(\T^2)\). By the Heinz inequality for self-adjoint operators (\cite[Proposition 10.14]{schmudgen_operators}) form equivalence of $L_c$ and $\Lambda$ translates to the form equivalence of $L_c^{s}$ and $\Lambda^{s}$ for $s \in [0,1]$. Then, Corollary 10.12 in \cite{schmudgen_operators} implies that $L_c^{-s}$ and $\Lambda^{-s}$ are form equivalent. In particular, for \(g\in L^2(\T^2)\),
\[
 \|g\|_{H^{-1+\sigma}}
 \asymp\|L_c^{-(1-\sigma)/2}g\|_{L^2}.
\]
For \(f\in\mathscr E_L\), \eqref{eq:Lc_sobolev_isomorphism} now gives
\[
 \|f\|_{H^{1+\sigma}}
 \asymp\|L_c f\|_{H^{-1+\sigma}}
 \asymp\|L_c^{(1+\sigma)/2}f\|_{L^2}.
\]
By Lemma~\ref{lem:eigenfunction_regularity}, \(\mathscr E_L\) is dense
in \(H^{1+\sigma}(\T^2)\), and it is dense in \(H_{L_c}^{1+\sigma}\) by spectral
truncation.  Both spaces embed continuously into \(L^2(\T^2)\), so completion identifies
them as the same space. Interpolation and duality give the remaining indices.
The assertion about Borel sets follows from equivalence of norms.
\end{proof}
%%%%

\subsection{Weak solution of the frozen SDE and its invariant measure}\label{sec_invdynamics_unique}
Throughout this subsection, all Hilbert spaces are real. We work on
\(\H:=\H_L\), equipped with the inner product
\eqref{eq:inner_product}, and set \(U:=L^2(\T^2)\).
By Corollary~\ref{cor:equ_norms}, \(\H\) is the original Sobolev space \(\H^{-\sigma}\) with an equivalent norm.
Define the operators 
\begin{equation}\label{eq:generator}
 \begin{split}
 \mathcal A(u,v)&=(v,-\eps^{-2}(Lu+v)),\\
 B_\eps h&=(0,\sqrt2\,\eps^{-2}h)
 \end{split}
\end{equation}
with domains $D(\mathcal A)=H_{L_c}^{1-\sigma}\times H_{L_c}^{-\sigma} =H^{1-\sigma}(\T^2)\times H^{-\sigma}(\T^2)$ and $D(B) = U$.
Here \(L=L_c-cI\) is extended continuously from
\(H_{L_c}^{1-\sigma}\) to \(H_{L_c}^{-1-\sigma}\).
This extension agrees with the form realization on their common
domain and does not require a definition of \(Mu\) for arbitrary
\(u\in H^{-\sigma}\).

We establish that~\eqref{eq:frozen_linear_dynamics} admits a weak solution in the sense of \cite{DaPrato_Zabczyk_2014} Sections 5.1--5.2. More precisely, on a filtered probability space, with a cylindrical \(U\)-Wiener process \(W\) relative
to the filtration, an adapted progressively measurable
\(\H\)-valued process \(X\) is a weak solution if its paths are locally
integrable in \(\H\) and, for every \(h\in D(\mathcal A^*)\),
\begin{equation}\label{eq:weak_frozen}
 \langle X(t),h\rangle_{\mathcal{H}}
 =\langle X(0),h\rangle_{\mathcal H}
  +\int_0^t\langle X(s),\mathcal A^*h\rangle_{\mathcal{H}}\,ds
  +\int_0^t\langle B_\eps^*h,dW_s\rangle_U .
\end{equation}
The identity is understood almost surely for each \(t\). We use
the continuous version obtained below. All adjoints in this
subsection are taken with respect to the fixed inner product of
\(\H\), unless explicitly stated otherwise.

%%%%GPT streamlined proof
\begin{proposition}[Existence and uniqueness of a weak solution]\label{prop:mild}
Suppose $M\in H^{-\sigma}(\T^2)$ for some $\sigma\in(0,\frac12)$.
The operator \(\mathcal A\) generates a strongly continuous semigroup \(S(t)\) for \(t\geq0\) on \(\H\).
For every \(\H\)-valued \(\mathcal F_0\)-measurable random variable
\(X_0\), the frozen equation has a unique weak solution, up to
modification, given by
\begin{equation}\label{eq:mild}
 X(t)=S(t)X_0+\int_0^t S(t-s)B_\eps\,dW_s
\end{equation}
for $t \geq 0$.

\noindent
%It has an adapted continuous \(\H\)-valued version.
\end{proposition}

\begin{proof} 
We show that the stochastic process $X(t)$ satisfies the conditions of Theorem 5.4 in \cite{DaPrato_Zabczyk_2014}. Then existence and uniqueness of a weak solution follow from this theorem. We are left to establish that $\mathcal A$ generates a strongly continuous semigroup in $\H$ and that the stochastic convolution is well-defined.
Decompose \(\mathcal A=\mathcal J+\mathcal R\), where
\[
 \mathcal J(u,v)=(v,-\eps^{-2}L_c u),\qquad
 \mathcal R(u,v)=(0,\eps^{-2}cu-\eps^{-2}v).
\]
On finite spectral sums, define, for every \(t\in\R\),
\begin{equation}\label{eq:abstract_wave_group}
 G(t)=
 \begin{pmatrix}
 \cos(tL_c^{1/2}/\eps)&
       \eps L_c^{-1/2}\sin(tL_c^{1/2}/\eps)\\
 -\eps^{-1}L_c^{1/2}\sin(tL_c^{1/2}/\eps)&
       \cos(tL_c^{1/2}/\eps)
 \end{pmatrix}.
\end{equation}
The spectral calculus interprets the off-diagonal entries between
the two components of \(\H\). On the \(k\)-th mode this is the
solution operator for
\(u_k'=v_k,\ v_k'=-\eps^{-2}\nu_k u_k\), where
\(\nu_k=\lambda_k+c\) are the eigenvalues of $L_c$. It preserves the equivalent norm
\begin{equation}\label{eq:equ_norm}
 \|(u,v)\|_{\H,\eps}^2
 =\sum_{k\geq1}\left(\nu_k^{-\sigma}|u_k|^2
             +\eps^2\nu_k^{-1-\sigma}|v_k|^2\right).
\end{equation}
Consequently, \(G(t)\) extends to an isometry for the norm
\(\|\cdot\|_{\H,\eps}\).
The trigonometric addition formulas give \(G(t+s)=G(t)G(s)\)
and \(G(-t)=G(t)^{-1}\). Strong continuity follows by approximation
with finite spectral sums.

Its generator is \(\mathcal J\). Indeed, differentiation in each
mode gives the stated action, and the condition
\(\mathcal J(u,v)\in \H\) is exactly
\[
 \sum_k\nu_k^{1-\sigma}|u_k|^2
       +\sum_k\nu_k^{-\sigma}|v_k|^2<\infty.
\]
For elements satisfying this condition, the spectral difference
quotients converge in \(\H\), by dominated convergence; conversely,
convergence in \(\H\) forces this condition by coordinatewise
differentiation. Thus, \(D(\mathcal J)=D(\mathcal A)\).

The embedding \(H_{L_c}^{-\sigma}\hookrightarrow H_{L_c}^{-1-\sigma}\)
makes \(\mathcal R\) bounded on \(\H\). Applying the bounded
perturbation theorem to \(\mathcal J+\mathcal R\) shows that \(\mathcal A\) generates a strongly continuous semigroups.
% Applying the bounded
% perturbation theorem to both \(\mathcal J+\mathcal R\) and
% \(-\mathcal J-\mathcal R\) shows that both \(\mathcal A\) and
% \(-\mathcal A\) generate strongly continuous semigroups, hence
% \(\mathcal A\) generates a strongly continuous group
% \cite[Theorems~III.1.3 and~II.3.11]{EngelNagel2000}.

To define the stochastic convolution as an element of
$L^2(\Omega;\mathcal H)$, it suffices to verify that
\[
    \int_0^t
    \|S(t-s)B_\eps\|_{\mathscr L_2(U,\mathcal H)}^2\,ds
    <\infty
    \qquad\text{for every }t>0.
\]
Here $\mathscr L_2(U,\mathcal H)$ denotes the space of
Hilbert--Schmidt operators from $U$ to $\mathcal H$, with
\[
    \|T\|_{\mathscr L_2(U,\mathcal H)}^2
    :=\sum_{k\geq1}\|Te_k\|_{\mathcal H}^2,
\]
where $\{e_k\}_{k \in \N}$ is any orthonormal basis of $U$.
The Hilbert space Itô isometry then gives
\[
    \mathbb E\left\|
        \int_0^t S(t-s)B_\eps\,dW_s
    \right\|_{\mathcal H}^2
    =
    \int_0^t
    \|S(t-s)B_\eps\|_{\mathscr L_2(U,\mathcal H)}^2\,ds.
\]
For that we return to the original inner
product \eqref{eq:inner_product}, without the additional factor
\(\eps^2\) in \eqref{eq:equ_norm}. In this norm,
\[
 \|B_\eps\|_{\mathscr L_2(U,\H)}^2
 =2\eps^{-4}\sum_{k\geq1}\nu_k^{-1-\sigma}<\infty,
\]
by Lemma~\ref{lemma_equ_ev} and \(\sum_{n\in\Z^2}(1+|n|^2)^{-1-\sigma}<\infty\).
Therefore, for every \(t>0\),
\[
 \int_0^t\|S(t-s)B_\eps\|_{\mathscr L_2(U,\H)}^2\,ds
 \leq t\sup_{0\leq r\leq t}\|S(r)\|_{\mathscr L(\H)}^2
             \|B_\eps\|_{\mathscr L_2(U,\H)}^2<\infty.
\]
This finishes the proof.
% For path continuity, the group property yields
% \[
%  \int_0^t S(t-s)B_\eps\,dW_s
%    =S(t)\int_0^t S(-s)B_\eps\,dW_s.
% \]
% The integral on the right has a continuous \(E\)-valued martingale
% version, because \(S(-s)B_\eps\) is locally square integrable in
% Hilbert--Schmidt norm. Strong continuity and local boundedness of
% \(S\) give continuity of the product. The deterministic term
% \(S(t)X_0\) is continuous almost surely as well.
\end{proof}

\begin{remark}
The shift \(c\) is used to construct equivalent norms, not to change
the drift. The operator in \eqref{eq:generator} contains \(L\),
not \(L+c\). Different admissible shifts yield the same realization:
they have equivalent Hilbert norms and the same action on the
finite spectral core. In particular, this construction does not
presuppose the positivity of \(L\) that is later deduced from
stationarity.
\end{remark}

% \begin{proposition}\label{prop_characterization_invmsr}
%     Let $\sigma \in (0, 1)$ and $\nu \in \mathcal{P}(\mathcal{H}_L)$. Assume $M_\mu \in H^{-\sigma}(\mathbb{T}^2)$.
%     Suppose $X(0)\sim \nu$ and that $\nu$ is an invariant measure for the flow 
%     \begin{equation*}
%         X(t) = e^{tA}X(0) + \int_0^t e^{(t-s)A} B \, dW_s,
%     \end{equation*}
%     which is well-defined by Proposition \ref{prop_invdynm_existence}.
%     Specifically, we assume $\operatorname{Law}(X(t)) = \nu$ for all $t \ge 0$. Then the following hold:
%     \begin{enumerate}
%         \item $L$ is strictly positive.
%         \item $\|e^{tA}\|_{\mathcal{H}_L} \le M e^{-\omega t}$ for some $\omega > 0, M \ge 1$, and $\nu$ is the unique invariant measure for the flow $X(t)$.
%         \item $\nu = \mathcal{N}\left(0, \begin{pmatrix} L^{-1} & 0 \\ 0 & \eps^{-2} \end{pmatrix}\right)$.
%     \end{enumerate}
% \end{proposition}

% \begin{proof}
%     From Theorem 11.7 (iii) in \cite{DaPrato_Zabczyk_2014}, the existence of an invariant measure implies that $\sup_{t \ge 0}\operatorname{Tr}(Q_t) < +\infty$. We use this to show that $L$ must be strictly positive. Suppose toward contradiction that $L$ is not strictly positive. Arguing as in Proposition \ref{prop_invdynm_existence}, we see that $L$ has compact resolvent and is bounded from below, hence, there exists a lowest eigenvalue $\mu_0 \le 0$ with a corresponding eigenfunction $e_0$. Recall the explicit solution for $e^{sA}B e_0$ discussed in Proposition \ref{prop_invdynm_existence}. Note, however, that here $e_0$ is an eigenfunction of $L$ and not of $L+c$ for some $c>0$ in contrast to Proposition \ref{prop_invdynm_existence}. If $\mu_0 \le 0$ then $\omega_0 = \frac{\sqrt{1 - 4 \eps^2 \mu_0}}{2 \eps^2} \geq \frac{1}{2\eps^2}$ is positive. Expanding the hyperbolic sine, we see that the solution $c_0(t)$ does not decay in $t$. It either approaches a non zero constant (if $\mu_0 = 0$) or exponentially (if $\mu_0 < 0$). In either case, the time integral diverges:
%     \begin{equation*}
%         \int_0^t \| e^{sA}B e_0 \|_{\mathcal{H}_L}^2 \, ds \to \infty \quad \text{as } t \to \infty.
%     \end{equation*}
%     Because $\int_0^t \| e^{sA}B e_0 \|_{\mathcal{H}_L}^2 \, ds \leq \operatorname{Tr}(Q_t)$, this contradicts the uniform bound $\sup_{t \ge 0}\operatorname{Tr}(Q_t) < \infty$.

%     It follows that $L$ is strictly positive. Consequently, as established in Proposition \ref{prop_invdynm_existence}, the real parts of the characteristic roots for all modes $m \in \mathbb{N}$ are strictly negative. Consequently, from computations similar to \eqref{esA_computation} and \eqref{bound_first_term}, we can infer that there exist constants $M \ge 1$ and $\omega > 0$ such that the semigroup decays exponentially:
%     \begin{equation}\label{etA_exponential_decay}
%         \| e^{tA} u \|_{\mathcal{H}_L} \le M e^{-\omega t} \|u\|_{\mathcal{H}_L} \quad \text{for all } t \ge 0 \text{ and } u \in \mathcal{H_L}.
%     \end{equation}
%     From Proposition 11.10 in \cite{DaPrato_Zabczyk_2014}, this strict exponential decay guarantees that $\nu$ is the unique invariant measure for the flow $X(t)$. 

%     Furthermore, Theorem 11.7 in \cite{DaPrato_Zabczyk_2014} establishes that this unique invariant measure takes the form $\mathcal{N}(0, \overline{Q})$, where the covariance operator $\overline{Q}$ is given by 
%     \begin{equation*}
%         \overline{Q} u \coloneq \lim_{t \to \infty} \int_0^t e^{sA}B B^* (e^{sA})^* u \, ds \quad \text{for } u \in \mathcal{H}_L.
%     \end{equation*}

%     We finish the proof by computing $\overline{Q}$ explicitly. Let $\{e_n\}_{n \in \N}$ be an orthonormal basis of $L^2(\T^2)$ of eigenfunctions of $L$ such that $L e_n = \mu_n$. We define the $L$-invariant 2-dimensional subspaces $V_n \coloneqq \operatorname{span}\{ (e_n, 0)^T, (0, e_n)^T \}$ of $\H_L$. Let $\overline{Q}_n \coloneqq \overline{Q}|_{V_n}$. On this subspace, $A$ restricts to the $2 \times 2$ matrix
%     \begin{equation*}
%         A_n \coloneqq A|_{V_n} = \begin{pmatrix} 0 & 1 \\ -\eps^{-2} \mu_n & -\eps^{-2} \end{pmatrix}.
%     \end{equation*}
%     By the construction of $e^{tA}$ and the $L$-invariance of $V_n$ it follows that $e^{tA}|_{V_n} = e^{tA_n}$. Furthermore, the basis elements of $V_n$ remain orthogonal in $\H_L$. Consequently, the union over $n \in \N$ of $V_n$ spans $\H_L$.

%     \noi
%     By the properties of the strongly continuous semigroup $e^{tA}$, we can differentiate the matrix exponential to obtain
%     \begin{equation*}
%         \frac{d}{dt} \left( e^{t A_n}B B^* e^{t A^*_n} \right) = A_n e^{t A_n}B B^* e^{t A^*_n} + e^{t A_n}B B^* e^{t A^*_n}A^*_n.
%     \end{equation*}
%     Integrating both sides with respect to $t$ over $[0, \infty)$ yields
%     \begin{equation*}
%         \int_0^\infty \frac{d}{dt} \left( e^{t A_n}B B^* e^{t A^*_n} \right) dt = \int_0^\infty \left( A_n e^{t A_n}B B^* e^{t A^*_n} + e^{t A_n}B B^* e^{t A^*_n}A^*_n \right) dt.
%     \end{equation*}
%     By the Fundamental Theorem of Calculus and since due to the exponential decay derived in \eqref{etA_exponential_decay}, the boundary term at infinity vanishes, the left side evaluates to
%     \begin{equation*}
%         \left[ e^{t A_n}B B^* e^{t A^*_n} \right]_0^\infty = 0 - B B^* = -B B^*.
%     \end{equation*}
%     Because the matrices $A_n$ and $A_n^*$ are independent of time and recalling the definition of $\overline{Q}_n$, we obtain 
%     \begin{align}\label{lyapunov_eq}
%         \begin{split}
%         -B B^* &= A_n \left( \int_0^\infty e^{t A_n}B B^* e^{t A^*_n} dt \right) + \left( \int_0^\infty e^{t A_n}B B^* e^{t A^*_n} dt \right) A^*_n \\
%         &= A_n \overline{Q}_n + \overline{Q}_n A_n^*.
%         \end{split}
%     \end{align}
%     Direct computations yield the adjoints on $\H_L$: 
%     \begin{align*}
%         A^*_n = \begin{pmatrix}
%             0 & 1 \\
%             -\eps^{-2} \mu_n & - \eps^{-2}
%         \end{pmatrix}, \quad \text{and} \quad
%         B^*_n =  \sqrt{2}\begin{pmatrix}
%             0 & \eps^{-2} \mu_n^{-1-\sigma}
%         \end{pmatrix}
%     \end{align*}
%     and the unique solution to this $2 \times 2$-matrix equation \eqref{lyapunov_eq} is
%     \begin{equation*}
%         \overline{Q}_n = \begin{pmatrix} \mu_n^{-1-\sigma} & 0 \\ 0 & \eps^{-2} \mu_n^{-1-\sigma}\end{pmatrix} = \begin{pmatrix} L^{-1-\sigma} & 0 \\ 0 & \eps^{-2} L^{-1-\sigma}\end{pmatrix}\Bigg|_{V_n}.
%     \end{equation*}
%     Since this holds for every subspace $V_n$ which together span $\mathcal{H}_L$, we conclude that 
%     $$\overline{Q} = \begin{pmatrix} L^{-1-\sigma} & 0 \\ 0 & \eps^{-2} L^{-1-\sigma}\end{pmatrix}.$$
%     Considering the definition of the inner product \eqref{eq:inner_product}, we can immediately deduce the covariance form, evaluated on $L^2(\T^2) \times L^2(\T^2)$ test functions, is given by $\begin{pmatrix} L^{-1} & 0 \\ 0 & \eps^{-2} \end{pmatrix}$. Consequently, $\nu$ is the stated Gaussian measure. 
% \end{proof}

%%%GPT Astra high rework of the previous Proposition 
\begin{proposition}[Characterization of the invariant measure]\label{prop_characterization_invmsr}
Fix \(0<\sigma<\frac12\), \(\eps>0\), and a real
\(M\in H^{-\sigma}(\T^2)\), and let \(L=\Lambda+M\) be its form
realization. Suppose \eqref{eq:frozen_linear_dynamics} has an invariant probability measure
\(\eta\in\mathcal P(\H^{-\sigma})\). Then:
\begin{enumerate}
 \item \emph{Positivity:} \(L\geq mI\) for some \(m>0\), and the forms of \(L\) and
       \(\Lambda\) are equivalent on \(H^1(\T^2)\).
 \item \emph{Exponential decay:} There exist \(C\geq1\) and \(\omega>0\) such that the operator norm in $\H$ of the semigroup $S(t)$ has exponential decay:
 \begin{equation}\label{etA_exponential_decay}
    \|S(t)\|_{\H}\leq Ce^{-\omega t},
       \qquad t\geq0.
 \end{equation}
 \item \emph{Gaussianity and uniqueness.} The invariant measure is unique and centered Gaussian.
       Its distributional covariance is
 \begin{equation}\label{eq:characterized_covariance}
       \begin{pmatrix}L^{-1}&0\\0&\eps^{-2}I\end{pmatrix}.
 \end{equation}
 In particular, for real \(f\in H^\sigma(\T^2)\) and
 \(g\in H^{1+\sigma}(\T^2)\),
 \begin{equation}\label{eq:characteristic_functional_mu}
  \int_{\mathcal H} e^{i\langle u,f\rangle+i\langle v,g\rangle}\,d\eta(u,v)
   =\exp\!\left[-\frac12\left(
       \langle L^{-1}f,f\rangle_{L^2_x}
          +\eps^{-2}\|g\|_{L^2_x}^2\right)\right]
 \end{equation}
 where the term on the left is understood by Sobolev duality.
\end{enumerate}
\end{proposition}

\begin{proof}
Set
\[
 Q_t=\int_0^t S(r)B_\eps B_\eps^*S(r)^*\,dr.
\]

\emph{Positivity of \(L\):}
For the positivity argument, choose $c\geq0$ such that
$L_c=L+cI$ is strictly positive, and equip $\H=\H_L$
with the inner product \eqref{eq:inner_product}.
Adjoints and traces in this step refer to this inner product.
The characterization theorem for invariant measures of linear
stochastic equations Theorem 11.7 in \cite{DaPrato_Zabczyk_2014}
implies, from existence of an invariant probability measure, that
\begin{equation}\label{eq:uniform_trace_bound}
 \sup_{t\geq0}\operatorname{Tr}_{\mathcal H} Q_t<\infty.
\end{equation}
Let $\{e_k\}_{k \in \N}$ be the spectral basis of $L$ with corresponding eigenvalues $\{\lambda_k\}_{k \in \N}$ then by the Hilbert--Schmidt trace identity,
\[
 \operatorname{Tr}_{\mathcal H} Q_t
   =\sum_{k\in\N}\int_0^t\|S(r)B_\eps e_k\|_{\mathcal H}^2\,dr.
\]
Suppose \(\lambda_k\leq0\). Write
\(S(t)B_\eps e_k=(a_k(t)e_k,a_k'(t)e_k)\). Then
\[
 \eps^2a_k''+a_k'+\lambda_ka_k=0,\qquad
 a_k(0)=0,\quad a_k'(0)=\sqrt2\,\eps^{-2}.
\]
With
\[
 r_\pm=\frac{-1\pm\sqrt{1-4\eps^2\lambda_k}}{2\eps^2},
\]
these distinct real roots give
\[
 a_k(t)=\frac{\sqrt2}{\eps^2(r_+-r_-)}
              \left(e^{r_+t}-e^{r_-t}\right).
\]
If \(\lambda_k=0\), this equals
\(\sqrt2(1-e^{-t/\eps^2})\); if \(\lambda_k<0\), it grows
exponentially. In either case,
\[
 \operatorname{Tr}_{\mathcal{H}} Q_t
 \geq (\lambda_k + c)^{-\sigma}\int_0^t|a_k(r)|^2\,dr
 \longrightarrow\infty,
\]
contrary to \eqref{eq:uniform_trace_bound}.
Thus, every \(\lambda_k>0\). Since \(L\) is bounded below and has
compact resolvent, \(m:=\inf_k\lambda_k>0\).
Moreover,
\[
 \frac{m}{m+c}L_c\leq L\leq L_c
\]
in the form sense. Lemma~\ref{lemma_equ_ev} proves the asserted
form equivalence.

\emph{Exponential decay:}
Since $L\geq mI$, we can choose $c=0$. 
To prove exponential decay of the semigroup we consider the $\eps$-adjusted energy functional
\[
 \mathcal E(u,v)
   :=\|L^{-\sigma/2}u\|_{L^2_x}^2
        +\eps^2\|L^{-(1+\sigma)/2}v\|_{L^2_x}^2
\]
that is a squared norm equivalent to \eqref{eq:inner_product} on $\H$. 

We first establish the energy estimate on finite spectral
subspaces. For $N\geq1$, set
\[
 \H_N:=\operatorname{span}\{(e_k,0),(0,e_k):1\leq k\leq N\}.
\]
The block-diagonal spectral representation of the generator (see the proof of Proposition~\ref{prop:mild})
implies that $\H_N$ is invariant under $S(t)$. Hence, for
$X_0\in \H_N$, the deterministic solution
\[
 X(t)=S(t)X_0=(u(t),v(t))
\]
has the finite expansions
\[
 u(t)=\sum_{k=1}^N u_k(t)e_k,
 \qquad
 v(t)=\sum_{k=1}^N v_k(t)e_k,
\]
where
\[
 u_k'(t)=v_k(t),\qquad
 \eps^2v_k'(t)=-\lambda_k u_k(t)-v_k(t).
\]
We refer to such a solution as a finite spectral solution.
All subsequent differentiations are therefore ordinary
differentiations of finite sums. Crucially, the estimates
below are independent of $N$ and of $X_0$.

Put
\[
 w(t)=L^{-(1+\sigma)/2}u(t),\qquad
 z(t)=L^{-(1+\sigma)/2}v(t).
\]
Then
\[
 w'(t)=z(t),\qquad \eps^2z'(t)=-Lw(t)-z(t),
\]
and the energy of the finite spectral solution is
\[
 \mathcal E(X(t)) 
 =\|L^{1/2}w(t)\|_{L^2_x}^2
    +\eps^2\|z(t)\|_{L^2_x}^2.
\]
For \(\delta>0\), define
\[
 \mathcal V(t)
   :=\mathcal E (X(t)) +2\delta\eps^2\langle w(t),z(t)\rangle_{L^2_x}
                    +\delta\|w(t)\|_{L^2_x}^2.
\]
Choose \(\delta\eps^2\leq\frac12\) and
\(\delta\eps/\sqrt m\leq\frac12\).
Cauchy--Schwarz and \(L\geq mI\) give
\[
 \frac12\mathcal E(X(t))\leq\mathcal V(t)
       \leq\left(\frac32+\frac{\delta}{m}\right)\mathcal E(X(t)).
\]
Direct differentiation yields
\[
 \mathcal V'(t)
  =-2(1-\delta\eps^2)\|z(t)\|_{L^2_x}^2
      -2\delta\|L^{1/2}w(t)\|_{L^2_x}^2
  \leq-\min\{\eps^{-2},2\delta\}\mathcal E(X(t)).
\]

Set
\[
 K:=\frac32+\frac{\delta}{m},
 \qquad q:=\min\{\eps^{-2},2\delta\}>0.
\]
The preceding estimates imply
\[
 \mathcal V'(t)
 \leq-q\mathcal E(X(t))
 \leq-\frac{q}{K}\mathcal V(t).
\]
Gronwall's inequality therefore yields
\[
 \mathcal E(S(t)X_0)
 \leq 2\mathcal V(t)
 \leq 2K e^{-qt/K}\mathcal E(X_0).
\]
By equivalence of norms,
\[
 \|S(t)X_0\|_{\mathcal H}
 \les \sqrt{2K} e^{-qt/(2K)}\|X_0\|_{\mathcal H}
\]
for $X_0\in\bigcup_{N\geq1}\H_N$.

\noindent
The constants are independent of the spectral cutoff.
Since $\bigcup_{N\geq1}\H_N$ is dense in $\H$ and $S(t)$
is bounded for each fixed $t\geq0$, this estimate extends
to every $X_0\in\H$. Taking the supremum over
$\|X_0\|_{\mathcal H}=1$ proves
\[
 \|S(t)\|_{\mathcal H}
 \leq C e^{-\omega t},
\]
for some $C\geq1$ and $\omega>0$.


\emph{Gaussianity and uniqueness:}
Utilizing the exponential decay we can apply Proposition 11.10 in \cite{DaPrato_Zabczyk_2014} to obtain that $\eta$ is the unique invariant measure of \eqref{eq:frozen_linear_dynamics}. Then, Theorem 11.7 in \cite{DaPrato_Zabczyk_2014} implies that $\eta = \mathcal{N}(0, Q_\infty)$ with $Q_\infty = \lim_{t\to\infty} Q_t$.
Finally, it remains to explicitly characterize $Q_\infty$. 
The vectors
\[
 F_k=(\lambda_k^{\sigma/2}e_k,0),\qquad
 G_k=(0,\lambda_k^{(1+\sigma)/2}e_k)
\]
form an orthonormal basis of \(\H\). The orthogonal planes
\(\operatorname{span}\{F_k,G_k\}\) reduce both \(S(t)\) and
\(B_\eps B_\eps^*\). On the \(k\)-th plane, the drift and diffusion
matrices in this orthonormal basis are
\[
 \widetilde A_k=
 \begin{pmatrix}
 0&\lambda_k^{1/2}\\
 -\eps^{-2}\lambda_k^{1/2}&-\eps^{-2}
 \end{pmatrix},\qquad
 \widetilde B_k=
 \begin{pmatrix}0\\ \sqrt2\,\eps^{-2}\lambda_k^{-(1+\sigma)/2}\end{pmatrix}.
\]
Thus, the covariance block \(C_k\) solves the usual Euclidean
Lyapunov equation
\begin{equation}\label{lyapunov_eq}
 \widetilde A_kC_k+C_k\widetilde A_k^{\mathsf T}
       +\widetilde B_k\widetilde B_k^{\mathsf T}=0.
\end{equation}
A direct substitution gives
\[
 C_k=\begin{pmatrix}
       \lambda_k^{-1-\sigma}&0\\
       0&\eps^{-2}\lambda_k^{-1-\sigma}
     \end{pmatrix}.
\]
There are no cross-mode covariance blocks, because the covariance
integrand is block diagonal. Consequently, the covariance operator
relative to the inner product of \(\H\) is
\begin{equation}\label{eq:hilbert_covariance}
 Q_\infty=
 \begin{pmatrix}
    L^{-1-\sigma}&0\\
    0&\eps^{-2}L^{-1-\sigma}
 \end{pmatrix}.
\end{equation}
For \(X=(u,v)\sim\mathcal N(0,Q_\infty)\), returning from the
orthonormal coordinates to the distributional coordinates yields
\[
 \mathbb E[\langle u,e_k\rangle\langle u,e_j\rangle]
     =\delta_{kj}\lambda_k^{-1},\qquad
 \mathbb E[\langle v,e_k\rangle\langle v,e_j\rangle]
     =\delta_{kj}\eps^{-2},
 \qquad
 \mathbb E[\langle u,e_k\rangle\langle v,e_j\rangle]=0.
\]
This proves \eqref{eq:characteristic_functional_mu} for finite
spectral test functions. Their density in
\(H^\sigma(\T^2)\times H^{1+\sigma}(\T^2)\), established in Lemma~\ref{lem:eigenfunction_regularity}, together with bounded convergence
on the left and continuity of the quadratic form on the right,
proves the identity for all stated test functions.
It is \eqref{eq:characterized_covariance}, rather than
\eqref{eq:hilbert_covariance}, that denotes the distributional
covariance of the invariant field.
\end{proof}


\subsection{Uniqueness by self-consistency}
Proposition \ref{prop_characterization_invmsr} reduces the remaining
argument to a self-consistency problem.  Namely, every admissible invariant law
has position marginal $\mathcal N(0,L^{-1})$. Utilizing this self-consistency relation we show that $M_\mu = 0$. This directly implies that $\Lambda = L$ which proves uniqueness. 
The following lemma identifies this
renormalized covariance diagonal and improves the regularity of
$M_\mu$ to $L^2(\T^2)$. 


%%%%GPT writing assistance 
\begin{lemma}\label{lem:renormalized_variance_convergence}
    Let $\sigma \in (0, \frac 12)$ and $u \sim \mathcal{N}(0, L^{-1})$ be a centered Gaussian random distribution. Assume that $M_\mu \in H^{-\sigma}(\T^2)$ and set $L = \Lambda + M_\mu$.

    Set
    \[
        \Gamma_N
        \coloneqq
        \{n\in\Z^2:\langle n\rangle\leq N\},
        \qquad
        \P_Nf
        \coloneqq
        \sum_{n\in\Gamma_N}
        \langle f,b_n\rangle_{L^2_x}b_n,
    \]
    where \(b_n(x)=e^{in\cdot x}\), and define
    \[
        M_N
        \coloneqq
        \mathbb E_\mu[(\P_Nu)^2]
        -
        \sum_{n\in\Gamma_N}\langle n\rangle^{-2}.
    \]
    Then $M_\mu\in L^2(\T^2)$ and
    $M_N\to M_\mu$ strongly in $L^2(\T^2)$.
    More precisely, writing $D=L^{-1}-\Lambda^{-1}$, we have
    \begin{equation}\label{eq:M_mu_expansion}
        M_\mu
        =
        \lim_{N\to\infty}
        \sum_{n\in\Gamma_N}\P_N(Db_n)b_{-n}
        =
        \lim_{N\to\infty}
        \sum_{n\in\Gamma_N}(Db_n)b_{-n},
    \end{equation}
    where both limits are taken in $L^2(\T^2)$.
\end{lemma}

\begin{proof}
    Since the covariance form of \(u\) is \(L^{-1}\), we have
    \[
        \mathbb E_\mu
        \left[
            \langle u,b_n\rangle_{L^2_x}
            \langle u,b_m\rangle_{L^2_x}
        \right]
        =
        \langle L^{-1}b_{-n},b_m\rangle_{L^2_x}.
    \]
    Consequently, after relabeling \(n\mapsto-n\),
    \[
        \mathbb E_\mu[(\P_Nu)^2]
        =
        \sum_{n\in\Gamma_N}
        \P_N(L^{-1}b_n)b_{-n}.
    \]
    Moreover,
    \[
        \sum_{n\in\Gamma_N}
        \P_N(\Lambda^{-1}b_n)b_{-n}
        =
        \sum_{n\in\Gamma_N}\langle n\rangle^{-2}.
    \]
    Hence,
    \begin{equation}\label{eq:renormalized_variance_D}
        M_N
        =
        \sum_{n\in\Gamma_N}
        \P_N(Db_n)b_{-n}.
    \end{equation}

    For \(m\in\Z^2\), define
    \[
        \Gamma_N(m)
        \coloneqq
        \{n\in\Z^2:
            n\in\Gamma_N,\ n+m\in\Gamma_N
        \}.
    \]
    Taking the \(m\)-th Fourier coefficient in
    \eqref{eq:renormalized_variance_D} gives
    \begin{equation}\label{eq:renormalized_variance_fourier}
        \widehat{M_N}(m)
        =
        \sum_{n\in\Gamma_N(m)}
        \langle Db_n,b_{n+m}\rangle_{L^2_x}.
    \end{equation}
    The resolvent identities, understood in the quadratic-form sense,
    yield
    \begin{align*}
        D
        &=
        -L^{-1}M_\mu\Lambda^{-1}
        \\
        &=
        -\Lambda^{-1}M_\mu\Lambda^{-1}
        +
        \Lambda^{-1}M_\mu
        L^{-1}M_\mu\Lambda^{-1}.
    \end{align*}
    Therefore,
    \[
        \widehat{M_N}(m)
        =
        -a_N(m)\widehat M_\mu(m)+R_N(m),
    \]
    where
    \[
        a_N(m)
        \coloneqq
        \sum_{n\in\Gamma_N(m)}
        \langle n\rangle^{-2}
        \langle n+m\rangle^{-2}
    \]
    and
    \begin{align*}
        R_N(m)
        \coloneqq
        \sum_{n\in\Gamma_N(m)}
        &\langle n\rangle^{-2}
        \langle n+m\rangle^{-2}
        \left\langle
            L^{-1}(M_\mu b_n),
            M_\mu b_{n+m}
        \right\rangle_{L^2_x}.
    \end{align*}

    We first consider the linear term. Define
    \[
        a(m)
        \coloneqq
        \sum_{n\in\Z^2}
        \langle n\rangle^{-2}
        \langle n+m\rangle^{-2}.
    \]
    Note, the following standard convolution estimate
    \begin{align}\label{eq:conv_ineq}
        \sum_{n \in \mathbb{Z}^2} \langle n \rangle^{-\alpha} \langle n + m \rangle^{-\alpha} \lesssim 
        \begin{cases}
             \jb{m}^{2-2\alpha} \quad \text{if $\alpha \in (1,2)$}
             \\
             \jb{m}^{-2} (1+\ln(\jb{m}^2)) \quad \text{if $\alpha = 2$.}
        \end{cases}
    \end{align}
    A proof is provided in Appendix Lemma \ref{lemma:conv_ineq}.
    This convolution estimate yields
    \begin{align}\label{eq:a_m_uniform_bound}
        \sup_{m\in\Z^2}a(m)<\infty.
    \end{align}
    For every fixed \(m\), we have
    \[
        a_N(m)\xlongrightarrow[]{N \to \infty} a(m),
        \qquad \text{and} \qquad
        0\leq a_N(m)\leq a(m).
    \]
    In particular, $\sup_m\langle m\rangle^\sigma a(m)<\infty$,
    so the assumption $M\in H^{-\sigma}(\T^2)$ implies
    \[
        \sum_{m\in\Z^2}
        a(m)^2|\widehat M_\mu(m)|^2
        \lesssim
        \|M_\mu\|_{H^{-\sigma}_x}^2<\infty.
    \]
    Dominated convergence therefore gives
    $a_N\widehat M\to a\widehat M$ in $\ell^2(\Z^2)$.


    It remains to treat \(R_N\). Multiplication by a Fourier mode
    satisfies
    \begin{equation}\label{eq:modulation_Hminus_sigma}
        \|M_\mu B_k\|_{H^{-\sigma}_x}
        \lesssim
        \langle k\rangle^\sigma
        \|M_\mu\|_{H^{-\sigma}_x}.
    \end{equation}
    Indeed, this follows directly from the standard algebraic inequality $\langle x+y \rangle^s \le 2^{-\frac{s}{2}} \langle x \rangle^s \langle y \rangle^{-s}$ for $s < 0$ after shifting Fourier frequencies.

    Next, we note that we obtain sufficient smoothing through the operator $L$: By Corollary~\ref{cor:equ_norms}, we have that $\Lambda^s$ and $L^s$ are form equivalent for all $s\in [-1,1]$.
    Furthermore, since $L^{-s}$ is a bounded for $s \geq 0$ and a self-adjoint operator on $L^2(\T^2)$, we have that $\langle L^{-s} u, u \rangle_{L^2_x} \les \|u\|^2_{L^2_x}$.
    Combined with \eqref{eq:modulation_Hminus_sigma}, the duality paring of $H^{\sigma}(\T^2)$ and $H^{-\sigma}(\T^2)$, these considerations yield:
    \begin{align*}
        \left|
        \left\langle
            L^{-1}(M_\mu b_n),
            M_\mu b_{n+m}
        \right\rangle
        \right|
        &\leq
        \|L^{-1}(M_\mu b_n)\|_{H^\sigma_x}
        \|M_\mu b_{n+m}\|_{H^{-\sigma}_x}
        \\
        &\lesssim
        \|M_\mu b_n\|_{H^{-1}_x}
        \|M_\mu b_{n+m}\|_{H^{-\sigma}_x}
        \\
        &\lesssim
        \langle n\rangle^\sigma
        \langle n+m\rangle^\sigma
        \|M_\mu\|_{H^{-\sigma}_x}^2.
    \end{align*}
    Hence, using the convolution estimate in \eqref{eq:conv_ineq} with $\alpha = 2-\sigma$, we have uniformly in \(N \in \N\),
    \begin{align}\label{eq:R_N_uniform_bound}
        |R_N(m)|
        &\lesssim
        \|M_\mu\|_{H^{-\sigma}_x}^2
        \sum_{n\in\Z^2}
        \langle n\rangle^{-2+\sigma}
        \langle n+m\rangle^{-2+\sigma}
        \nonumber\\
        &\lesssim
        \|M_\mu\|_{H^{-\sigma}_x}^2
        \langle m\rangle^{-2+2\sigma}.
    \end{align}
    Thus, the full series defining $R(m)$ is absolutely
    convergent for every $m$, and $R_N(m)\to R(m)$.
    Since $\sigma<\frac12$,
    \[
        \sum_{m\in\Z^2}\langle m\rangle^{-4+4\sigma}<\infty.
    \]
    Dominated convergence consequently gives
    $R_N\to R$ in $\ell^2(\Z^2)$.

    It follows by Parseval's identity that $M_N$ converges
    in $L^2(\T^2)$ to $M_\mu$. Hence, $M_\mu\in L^2(\T^2)$ and
    $M_N\to M_\mu$ strongly in $L^2(\T^2)$.


    Finally, for
    \[
        \widetilde M_N
        =
        \sum_{n\in\Gamma_N}(Db_n)b_{-n},
    \]
    the same Fourier computation applies with
    $\Gamma_N(m)$ replaced by $\Gamma_N$ in the definitions
    of $a_N(m)$ and $R_N(m)$.
    The pointwise limits remain $a(m)$ and $R(m)$, and
    all the preceding bounds remain valid.
    The same dominated-convergence argument therefore
    gives $\widetilde M_N\to M$ in $L^2(\T^2)$, proving
    both limits in \eqref{eq:M_mu_expansion}.

\end{proof}


Finally, we are prepared to prove the conditional uniqueness of the invariant measure of the SDMFW.
\begin{proof}[Proof of Theorem \ref{theorem:uniqueness_SDMFW}]
    Suppose there exists an invariant measure $\mu \in \mathcal{W}$ of the \eqref{eq:SDMFW}, and set $M_\mu=\mathbb E_\mu[\wick{u^2}]$ and $L=\Lambda+M_\mu$. 
    Due to the equivalence of Hilbert spaces established in Corollary~\ref{cor:equ_norms} we can apply Proposition~\ref{prop_characterization_invmsr} to infer that $L$ is strictly positive and that $\mu = \mathcal{N}(0, L^{-1}) \otimes \mathcal{N}(0, \varepsilon^{-2}I)$. 
    
    Applying Lemma \ref{lem:renormalized_variance_convergence}
    to the position marginal, we obtain
    $M_\mu\in L^2(\T^2)$ and
    \[
        \widetilde M_N
        \coloneqq
        \sum_{n\in\Gamma_N}
        (L^{-1}-\Lambda^{-1})b_n\,b_{-n}
        \longrightarrow M_\mu
        \qquad\text{in }L^2(\T^2).
    \]

    Since $M_\mu$ is real-valued, the resolvent identity
    $L^{-1}-\Lambda^{-1}=-L^{-1}M_\mu\Lambda^{-1}$
    gives, for every $N$,
    \begin{align*}
        \langle\widetilde M_N,M_\mu\rangle_{L^2_x}
        &=
        \sum_{n\in\Gamma_N}
        \left\langle
            (L^{-1}-\Lambda^{-1})b_n,M_\mu b_n
        \right\rangle_{L^2_x}\\
        &=
        -\sum_{n\in\Gamma_N}\langle n\rangle^{-2}
        \left\langle
            L^{-1}(M_\mu b_n),M_\mu b_n
        \right\rangle_{L^2_x}\\
        &\leq0.
    \end{align*}
    Here all sums are finite, $M_\mu b_n\in L^2(\T^2)$,
    and the last inequality follows from positivity of
    $L^{-1}$.
    Passing to the limit using the strong $L^2$ convergence
    of $\widetilde M_N$ yields
    \[
        \|M_\mu\|_{L^2_x}^2
        =
        \lim_{N\to\infty}
        \langle\widetilde M_N,M_\mu\rangle_{L^2_x}
        \leq0.
    \]
    Consequently, $M_\mu=0$, so $L=\Lambda$ and
    $\mu=\rho_\eps$.

    Conversely, the chosen Wick ordering gives
    $\mathbb E_{\rho_\eps}[\wick{u^2}]=0$.
    The stationary linear damped-wave dynamics with
    position covariance $\Lambda^{-1}$ and velocity
    covariance $\eps^{-2}I$ therefore provide a stationary
    solution of \eqref{eq:SDMFW} with law $\rho_\eps$.
    Hence, $\rho_\eps$ is the unique invariant measure
    in $\mathcal W$.
\end{proof}


%%%% Author's version of this section before GPT's rework 
% \subsection{Convergence of SDMFW's Invariant Measure to the SMFW's Invariant Measure}\label{sec_invdynamics_SKA}

% As discussed in Section \ref{Preperation_Invariant_Measure}, we show that the SKA also occurs for the respective probability measures.
% Let $\sigma > 0$ and $\pi: \mathcal{H}^{-\sigma}(\T^2) \to H^{-1-\sigma}(\T^2)$ be the canonical projection onto the second component.
% To prove convergence, we define a random variable $\psi_\eps \sim \mathcal{N}(0, \eps^{-2}I) = \pi_* \rho_\eps$ on $H^{-1-\sigma}(\T^2)$. Let $\{h_n\}_{n \in \Z^2}$ be a family of independent random variables such that for $n \neq 0$, $h_n$ is a complex-valued centered Gaussian random variable with variance $2$ satisfying $h_{-n} = \overline{h_n}$ for $n \in \mathcal{J}$, and for $n = 0$, $h_0$ is a real-valued centered Gaussian random variable with variance $1$. We assume these are independent of the cylindrical Wiener process $W_t$ defined in \eqref{Wiener1}.
% Define
% \begin{align*}
%     \psi_\eps(x) \coloneqq  \eps^{-1}h_0 + \frac{1}{\sqrt{2}}\sum_{n \in \Z^2 \setminus\{0\}} \eps^{-1} h_n e^{i n \cdot x} \sim \mathcal{N}(0, \eps^{-2}I).
% \end{align*}

% First, we show that $(\P_N \psi_\eps)_*\mathbb{P}$ converges weakly to $(\psi_\eps)_*\mathbb{P} = \pi_* \rho_\eps$ as $N \to \infty$, where $\P_N$ is the truncated Fourier projection defined in \eqref{P_N}. This weak convergence is implied by the following standard statement.

% \begin{lemma}\label{prop_measure_well_definedness}
%     For fixed $\eps \in (0,1]$, we have
%     \begin{align*}
%         \P_N \psi_\eps \rightarrow \psi_\eps \quad \text{in } L^2(\Omega; H^{-1-\sigma}(\T^2)),
%     \end{align*}
%     as $N \to \infty$.
% \end{lemma}
% \begin{proof}
%     For $s \geq 0$ and $N \in \N$, define the truncated random variable
%     \begin{align}\label{truncated_psi}
%         \psi_{N,s}(x) \coloneq \sum_{\substack{n \in \Z^2 \\ \jb{n} \leq N}} \frac{h_n}{\jb{n}^s} e^{ i n \cdot x},
%     \end{align}
%     where $h_n$ is defined as above. Let $\sigma \in \R$. We will show that $\psi_{N,s}$ is a Cauchy sequence in $L^2(\Omega; H^{\sigma}(\T^2))$. Let $M, N \in \N$ with $M \leq N$. Using the independence of $h_n$ and the fact that $\EV{|h_n|^2} \leq 2$, we compute the expectation of the difference
%     \begin{align*}
%         \EV{\|\psi_{N,s} - \psi_{M,s} \|^2_{H^\sigma_x}} 
%         &= 
%         \sum_{\substack{n \in \Z^2 \\ M < \jb{n} \leq N}} \jb{n}^{2\sigma} \EV{\Big| \frac{h_n}{\jb{n}^s} \Big|^2}
%         \\ &\leq
%         2 \sum_{\substack{n \in \Z^2 \\ M < \jb{n} \leq N}} \jb{n}^{2 \sigma - 2 s}.
%     \end{align*}
%     This sum converges if and only if $\sigma < s - 1$.
%     Hence, for any $\sigma <  s-1$, the limit $\lim_{N\to \infty}\psi_{N,s}$ is well-defined in $L^2(\Omega; H^{\sigma}(\T^2))$. Since our target space requires $-1-\sigma < -1$, the statement follows directly. 
% \end{proof}

% We are now ready to prove the convergence of the probability measures. 

% \begin{theorem}
%     Let $\sigma > 0$, $\eps > 0$, and $\pi: \mathcal{H}^{-\sigma}(\T^2) \to H^{-1-\sigma}(\T^2)$ be the canonical projection onto the second component. Let $\rho_\eps \coloneqq \mathcal{N}\left(0, \begin{pmatrix}\Lambda^{-1} & 0 \\ 0 & \eps^{-2}\end{pmatrix}\right)$ and $\rho_0 \coloneqq \mathcal{N}(0, \Lambda^{-1})$.
%     Then $\rho_\eps \rightarrow \rho_0$ as $\eps \to 0$ in the sense of Definition \ref{def_invdynamics_SKA}.
% \end{theorem}
% \begin{proof}
%     Let $f \in C_0(H^{-1-\sigma}(\T^2))$ be arbitrary. By the definition, $f$ is bounded, thus there exists $C>0$ such that $|f| \leq C$. Furthermore, for any $\delta > 0$, there exists a $K> 0$ such that $|f(u)| < \delta$ for all $\|u\|_{H^{-1-\sigma}_x} > K$. 
    
%     Using the fact that $(\pi_* \rho_\eps)$ is a probability measure, therefore it is bounded by $1$, a Chebyshev-like inequality, and the fact that $(\P_N \psi_\eps)_* \mathbb{P}$ converges weakly to $(\psi_\eps)_* \mathbb{P}$ (by Lemma \ref{prop_measure_well_definedness}), we obtain  
%     \begin{align*}
%     \begin{split}
%        \bigg| \int_{H^{-1-\sigma}(\T^2)} f(u) \, d(\pi_* \rho_\eps)(u) \bigg| 
%        &\leq \delta (\pi_* \rho_\eps)(\|u\|_{H^{-1-\sigma}_x} > K) +  C (\pi_* \rho_\eps)\big( \|u \|_{H^{-1-\sigma}_x} \leq K \big) 
%        \\
%        &= \delta + C(\pi_* \rho_\eps)\big( \|u \|_{H^{-1-\sigma}_x} \leq K \big)
%        \\
%         &\leq \delta + C e^{K^2} \int_{H^{-1-\sigma}(\T^2)} e^{-\|u \|^2_{H^{-1-\sigma}_x}} \, d(\psi_\eps)_*\mathbb{P}(u)
%         \\ &=
%         \delta + C e^{K^2} \lim_{N \to \infty}\int_{\Omega}e^{-\|\P_N \psi_\eps(\omega) \|^2_{H^{-1-\sigma}_x}} \, d \mathbb{P}(\omega).
%     \end{split}
%     \end{align*}
    
%     To compute the last term, we denote $\mathcal{J}_N \coloneq \{ n \in \mathcal{J} \mid \jb{n} \leq N \}$. This is a finite subset of $\Z^2$. Furthermore, by the cylindrical construction of the random variables $h_n$, namely that $h_{-n} = \overline{h_n}$ for $n \in \mathcal{J}$, we have $|h(n)|^2 = |h(-n)|^2$. Consequently, we obtain
%     \begin{align*}
%         \lim_{N \to \infty} &\int_{\Omega} e^{-\|\P_N \psi_\eps(\omega) \|^2_{H^{-1-\sigma}_x}} \, d \mathbb{P}(\omega)
%         \\
%         &= \lim_{N \to \infty} \int_\Omega \exp\left(-\jb{0}^{-2-2\sigma} \eps^{-2} |h_0(\omega)|^2 - \sum_{n \in \mathcal{J}_N} \jb{n}^{-2-2\sigma} \eps^{-2} |h_n(\omega)|^2 \right) \, d\mathbb{P}({\omega})
%         \\
%         &= \lim_{N \to \infty} (1+ 2\eps^{-2})^{-\frac{1}{2}} \prod_{n \in \mathcal{J}_N} (1+ 2\eps^{-2} \jb{n}^{-2-2\sigma})^{-1}
%         \\ 
%         &= (1+ 2\eps^{-2})^{-\frac{1}{2}} \lim_{N \to \infty} \prod_{n \in \mathcal{J}_N}\Big(1 - \frac{2}{\eps^2 \jb{n}^{2+2\sigma} +2}\Big) . 
%     \end{align*}

%     % We note two technical points regarding the previous estimate. The first inequality arises because the domain of integration on the left-hand side is restricted to $\|u\|_{H^{-1-\sigma}_x} < \infty$, while the right-hand side is an integral over the entire space. Second, the variable $x$ represents the coefficient vector in $\R^{|\Z^2_N|}$ indexed by $n \in \Z^2_N$.
%     \noi
%     This infinite product converges to a strictly positive number if and only if the series
%     \begin{align*}
%         \sum_{n \in \Z^2} \frac{2}{\eps^2 \jb{n}^{2+2\sigma} + 2} < \infty
%     \end{align*}
%     converges. This holds true because $2 + 2\sigma > 2$. Note, that the product converges for all $\eps \in [0,1]$, since for $\eps = 0$ the product equals $0$. Finally, as $(1+ 2\eps^{-2})^{-\frac{1}{2}} \to 0$ for $\eps \to 0$ and since the infinite product is uniformly bounded in $\eps$ the final expression converges to $0$. Consequently, the integral converges to $0$, implying that $ \rho_\eps \rightarrow \rho_0$ in the sense of Definition \ref{def_invdynamics_SKA}.
% \end{proof}


%%%%GPT Astra high streamlined the original proof (which can be found directly above)
\section{Small-mass behavior of the invariant measures}\label{sec_invdynamics_SKA}

Recall that
\[
    \rho_\eps
    =
    \mathcal{N}(0, \Lambda^{-1}) \otimes \mathcal{N}(0, \varepsilon^{-2}I),
    \qquad \text{and} \qquad
    \rho_0=\mathcal N(0,\Lambda^{-1}).
\]
The uniqueness of the invariant measure $\rho_0$ of the \eqref{eq:SMFH} follows analogous to the uniqueness of the invariant measure of the \eqref{eq:SDMFW}.
Concerning the Smoluchowski-Kramers type behavior of these invariant measures, the position marginal of \(\rho_\eps\) equals \(\rho_0\) for every
\(\eps>0\). By contrast, the velocity marginal escapes to infinity as \(\eps\to0\). This will be demonstrated rigorously in the following proof.

\begin{proof}[Proof of Theorem \ref{theorem:convergence_uniqueness_SMFH}]
%%%%GPT Astra High generalization of self-consistency argument to the SMFH
    We first establish uniqueness of the invariant measure of the
\eqref{eq:SMFH}. Since this proof follows the exact same steps presented in Section~\ref{Chpt_Invariant_Dynamics} we only discuss the important steps. 
Let $\mu\in\mathcal W_0$ be an invariant measure,
and set $M_\mu\coloneqq\mathbb E_\mu[\wick{u^2}]$ and
$L\coloneqq\Lambda+M_\mu$. By stationarity, the equation reduces to
\[
    du(t)=-Lu(t)\,dt+\sqrt{2}\,dW_t.
\]
Lemmas \ref{lem:relative_bound} and \ref{lemma_equ_ev} provide a
self-adjoint realization of $L$ with compact resolvent. If
$Le_j=\lambda_j e_j$ in a real orthonormal eigenbasis, the
corresponding coordinates satisfy
\[
    d\langle u(t),e_j\rangle
    =-\lambda_j\langle u(t),e_j\rangle\,dt+\sqrt{2}\,d\beta_j(t),
\]
where the $\beta_j$ are independent standard Brownian motions.
The existence of an invariant probability measure excludes
$\lambda_j\leq0$, since the corresponding scalar equation has
no invariant probability measure. Thus, $L$ is strictly positive.
Every finite collection of these coordinates consequently has
the unique invariant law given by the product of the centered
Gaussian measures with variances $\lambda_j^{-1}$. It follows
that $\mu=\mathcal N(0,L^{-1})$, with covariance understood in
the distributional sense. The remainder of the argument is
identical to the self-consistency argument in the proof of
Theorem \ref{theorem:uniqueness_SDMFW}: Lemma
\ref{lem:renormalized_variance_convergence} yields
$M_\mu\in L^2(\T^2)$, and the resolvent identity together with
positivity of $L^{-1}$ gives $M_\mu=0$. Hence, $L=\Lambda$ and
$\mu=\rho_0$. Conversely, the chosen Wick ordering gives
$\mathbb E_{\rho_0}[\wick{u^2}]=0$, so the stationary
Ornstein-Uhlenbeck process with drift $-\Lambda$ and covariance
$\Lambda^{-1}$ is a stationary solution of \eqref{eq:SMFH}.
Therefore, $\rho_0$ is the unique invariant measure in
$\mathcal W_0$.

%%%%%
Next, we consider the stated Smoluchowski-Kramers type behavior.
    The identity
    \[
        (\pi_1)_*\rho_\eps=\rho_0
    \]
    follows immediately from the definition of \(\rho_\eps\).

    It remains to consider the velocity marginal. Let
    \(\psi_1\) be a centered Gaussian random distribution whose
    covariance form on \(L^2(\T^2)\) is the identity. 
    Let $\{h_n\}_{n \in \Z^2}$ be a family of random variables that are independent for $n \in (\Z_+\times \{0\}) \cup (\Z\times \Z_+) \cup\{(0, 0)\}$ such that for $n \neq 0$, $h_n$ is a complex-valued centered Gaussian random variable with variance $2$ satisfying $h_{-n} = \overline{h_n}$, and for $n = 0$, $h_0$ is a real-valued centered Gaussian random variable with variance $1$.
    Using these Gaussian variables, it may be realized as
    \[
        \psi_1(x)
        =
        h_0
        +
        \frac{1}{\sqrt{2}}
        \sum_{n\in\Z^2\setminus\{0\}}
        h_ne^{in\cdot x}.
    \]
    Since \(\sigma>0\),
    \[
        \mathbb E\big[
            \|\psi_1\|_{H^{-1-\sigma}_x}^2
        \big]
        =
        \sum_{n\in\Z^2}
        \langle n\rangle^{-2-2\sigma}
        <\infty.
    \]
    Hence, the series defining \(\psi_1\) converges in
    \(L^2(\Omega;H^{-1-\sigma}(\T^2))\).

    For \(\eps>0\), define
    \[
        \psi_\eps\coloneqq\eps^{-1}\psi_1.
    \]
    Then \(\psi_\eps\) has covariance form \(\eps^{-2}I\), and
    therefore
    \[
        \operatorname{Law}(\psi_\eps)
        =
        (\pi_2)_*\rho_\eps.
    \]
    Moreover,
    \[
        \|\psi_\eps\|_{H^{-1-\sigma}_x}
        =
        \eps^{-1}
        \|\psi_1\|_{H^{-1-\sigma}_x}.
    \]
    The Gaussian field \(\psi_1\) is nonzero almost surely. Indeed,
    \(\psi_1=0\) would in particular imply \(h_0=0\), which has
    probability zero. Consequently,
    \[
        \|\psi_\eps\|_{H^{-1-\sigma}_x}
        \longrightarrow\infty
        \qquad\text{almost surely as }\eps\to0.
    \]

    Let
    \[
        f\in
        C_{\mathrm{rv}}
        (H^{-1-\sigma}(\T^2)).
    \]
    By the radial-vanishing property,
    \[
        f(\psi_\eps)\longrightarrow0
        \qquad\text{almost surely}.
    \]
    Since \(f\) is bounded, the dominated convergence theorem gives
    \begin{align*}
        \int_{H^{-1-\sigma}(\T^2)}
        f(v)\,d((\pi_2)_*\rho_\eps)(v)
        = \mathbb E[f(\psi_\eps)]
        \longrightarrow0.
    \end{align*}
    This proves the asserted convergence.
\end{proof}